\section{Appendix}


\subsection{Experimental Environment}
All experiments are conducted under a unified hardware and software configuration to ensure fair comparison across methods and benchmarks.
Unless otherwise stated, the evaluation environment is as follows.

\begin{itemize}
    \item \textbf{GPU:} NVIDIA RTX 5880 Ada GPU with 48\,GB memory.
    \item \textbf{CPU:} Dual-socket AMD EPYC 9654 processors (96 cores per socket, 192 cores total, 384 threads).
    \item \textbf{CUDA / Driver:} CUDA~13.0 with NVIDIA driver version~580.65.06.
\end{itemize}

All latency measurements are obtained using CUDA events and exclude environment simulation and robot actuation overhead.

\subsection{Vision Encoder Configuration}
Unless otherwise specified, the visual observation at each timestep is resized to $224\times224$ and processed by a ViT-based vision encoder.
Images are partitioned into non-overlapping patches of size $14\times14$, resulting in
\[
P = (224 / 14)^2 = 256
\]
visual patches per frame.
We evaluate two widely used vision backbones in OpenVLA-style models: DINOv2 and SigLIP.
These configurations are kept identical across baselines and CTR-enabled variants to isolate the effect of coordinated token reuse.

\subsection{CTR Configuration and Latency Accounting}
Unless otherwise specified, CTR uses grayscale patch difference for temporal consistency, a difference threshold of $0.004$, top-$K$ task-relevance selection with $K=160$, and a keyframe interval of $5$ control steps.
Patch-level attention scores are computed from the previous-step text-to-vision attention: attention maps are averaged over heads, text-query positions are averaged over the vision-token columns, and the top-$K$ visual patches are marked task-relevant and forced to recompute.
The LLM-side layer schedule follows the entropy-based cache-reuse schedule used by the VLA-cache component.

OpenVLA-style models use two vision-encoder branches, DINOv2 and SigLIP.
Accordingly, main-table model-core latency and TFLOPs are computed as
\[
    \mathrm{Total} = \mathrm{VE}_{\mathrm{DINOv2}} + \mathrm{VE}_{\mathrm{SigLIP}} + \mathrm{LLM}.
\]
When a component ablation reports a single VE latency or TFLOPs value, it is reported per vision branch; this is why the main-table totals correspond to $2\times$VE + LLM in the OpenVLA ablation.
All reported efficiency numbers exclude simulator stepping, rendering, and robot actuation.

\subsection{Benchmarks}
We evaluate CTR on two complementary simulation benchmarks that cover a wide spectrum of manipulation behaviors.

\textbf{LIBERO} is used as the primary benchmark and includes four task suites:
\textit{Spatial} (precise object placement and spatial reasoning),
\textit{Object} (single-object manipulation with diverse geometries),
\textit{Goal} (goal-conditioned, multi-step manipulation),
and \textit{Long} (long-horizon tasks requiring temporal consistency).
For LIBERO experiments, we use the official OpenVLA-7B checkpoint that has been fine-tuned on the corresponding LIBERO task suites.

To assess generalization beyond LIBERO, we additionally evaluate on \textbf{SimplerEnv}, a tabletop manipulation benchmark with different object layouts and visual statistics.
We report results on three representative tasks: \textit{Move Near Object}, \textit{Pick Coke Can}, and \textit{Drawer Operations}.
For SimplerEnv, we use the original OpenVLA-7B checkpoint without any task-specific fine-tuning.
All CTR hyperparameters are kept identical to those used in LIBERO experiments.

\subsection{Model Variants and Checkpoints}
We evaluate two VLA backbones: OpenVLA and OpenVLA-OFT.
All experiments use officially released OpenVLA-7B checkpoints.
CTR is applied strictly at inference time and does not modify model weights, training data, or optimization procedures.
Both the baseline and CTR-enabled models share identical architectures and token layouts.


\subsection{Qualitative Visualization of an Inference Trajectory}
Figure~\ref{fig:appendix_trajectory} visualizes a complete closed-loop VLA inference trajectory from a single manipulation episode. The figure highlights that most visual regions remain unchanged across consecutive timesteps, which motivates exploiting temporal redundancy via token reuse.

\begin{figure*}[t]
    \centering
    \includegraphics[width=1.0\linewidth]{figures/grid_2x6_phase.png}
    \caption{A complete closed-loop VLA inference trajectory from a single manipulation episode. Consecutive frames exhibit strong temporal redundancy, with only localized regions changing over time.}
    \label{fig:appendix_trajectory}
\end{figure*}
